\documentclass{amsart}
% For dealing with references we use the comment environment
\usepackage{verbatim}
\newenvironment{reference}{\comment}{\endcomment}
%\newenvironment{reference}{}{}
\newenvironment{slogan}{\comment}{\endcomment}
\newenvironment{history}{\comment}{\endcomment}

% For commutative diagrams we use Xy-pic
\usepackage[all]{xy}

% We use 2cell for 2-commutative diagrams.
\xyoption{2cell}
\UseAllTwocells

% We use multicol for the list of chapters between chapters
\usepackage{multicol}

% This is generally recommended for better output
\usepackage{lmodern}
\usepackage[T1]{fontenc}

% For cross-file-references

% Package for hypertext links:
\usepackage{hyperref}

% For any local file, say "hello.tex" you want to link to please

% Theorem environments.
%
\theoremstyle{plain}
\newtheorem{theorem}[subsection]{Theorem}
\newtheorem{proposition}[subsection]{Proposition}
\newtheorem{lemma}[subsection]{Lemma}

\theoremstyle{definition}
\newtheorem{definition}[subsection]{Definition}
\newtheorem{example}[subsection]{Example}
\newtheorem{exercise}[subsection]{Exercise}
\newtheorem{situation}[subsection]{Situation}

\theoremstyle{remark}
\newtheorem{remark}[subsection]{Remark}
\newtheorem{remarks}[subsection]{Remarks}

\numberwithin{equation}{subsection}

% Macros
%
\def\lim{\mathop{\mathrm{lim}}\nolimits}
\def\colim{\mathop{\mathrm{colim}}\nolimits}
\def\Spec{\mathop{\mathrm{Spec}}}
\def\Hom{\mathop{\mathrm{Hom}}\nolimits}
\def\Ext{\mathop{\mathrm{Ext}}\nolimits}
\def\SheafHom{\mathop{\mathcal{H}\!\mathit{om}}\nolimits}
\def\SheafExt{\mathop{\mathcal{E}\!\mathit{xt}}\nolimits}
\def\Sch{\mathit{Sch}}
\def\Mor{\mathop{\mathrm{Mor}}\nolimits}
\def\Ob{\mathop{\mathrm{Ob}}\nolimits}
\def\Sh{\mathop{\mathit{Sh}}\nolimits}
\def\NL{\mathop{N\!L}\nolimits}
\def\CH{\mathop{\mathrm{CH}}\nolimits}
\def\proetale{{pro\text{-}\acute{e}tale}}
\def\etale{{\acute{e}tale}}
\def\QCoh{\mathit{QCoh}}
\def\Ker{\mathop{\mathrm{Ker}}}
\def\Im{\mathop{\mathrm{Im}}}
\def\Coker{\mathop{\mathrm{Coker}}}
\def\Coim{\mathop{\mathrm{Coim}}}

% Boxtimes
%
\DeclareMathSymbol{\boxtimes}{\mathbin}{AMSa}{"02}

%
% Macros for moduli stacks/spaces
%
\def\QCohstack{\mathcal{QC}\!\mathit{oh}}
\def\Cohstack{\mathcal{C}\!\mathit{oh}}
\def\Spacesstack{\mathcal{S}\!\mathit{paces}}
\def\Quotfunctor{\mathrm{Quot}}
\def\Hilbfunctor{\mathrm{Hilb}}
\def\Curvesstack{\mathcal{C}\!\mathit{urves}}
\def\Polarizedstack{\mathcal{P}\!\mathit{olarized}}
\def\Complexesstack{\mathcal{C}\!\mathit{omplexes}}
% \Pic is the operator that assigns to X its picard group, usage \Pic(X)
% \Picardstack_{X/B} denotes the Picard stack of X over B
% \Picardfunctor_{X/B} denotes the Picard functor of X over B
\def\Pic{\mathop{\mathrm{Pic}}\nolimits}
\def\Picardstack{\mathcal{P}\!\mathit{ic}}
\def\Picardfunctor{\mathrm{Pic}}
\def\Deformationcategory{\mathcal{D}\!\mathit{ef}}

\setlength{\emergencystretch}{3em}
\begin{document}
\title[Stacks Project, Tag 0EB6]{312 --- Hartshorne-Lichtenbaum top local-cohomology vanishing}
\maketitle
\noindent Copyright (C) 2005--2025 Johan de Jong. Stacks Project authors. Source: \href{https://stacks.math.columbia.edu/tag/0EB6}{Tag 0EB6}.

\noindent Editorial difficulty note (not author text). Difficulty H3: The proof reduces through completion and maximal-dimensional components to a complete local domain. It then replaces the support ideal by a one-dimensional prime and changes ordinary powers to symbolic powers in a dualizing Ext colimit; the depth of these symbolic quotients forces vanishing. The task requires the duality-and-symbolic-power mechanism rather than ordinary cohomological dimension bounds..

\noindent Editorial provenance note (not author text). History: excerpt from revision \texttt{a0444\allowbreak 6e57e\allowbreak c1fbc\allowbreak 252a8\allowbreak 71afc\allowbreak ec775\allowbreak 2fb28\allowbreak 07b14}, local-cohomology.tex, lines 3504--3556. Reference presentation was adapted; original-author context and core support blocks were retained or added. The mathematical wording is unchanged. Licensed under GNU FDL 1.2 or later; no invariant sections or cover texts. See ../licenses/COPYING. Context blocks reproduce author definitions and supporting results; they are not additional counted problems.

\begin{lemma}
\label{lemma-cd-bound-dualizing}
Let $A$ be a Noetherian local ring with
normalized dualizing complex $\omega_A^\bullet$.
Let $I \subset A$ be an ideal.
If $H^0_{V(I)}(\omega_A^\bullet) = 0$, then $\text{cd}(A, I) < \dim(A)$.
\end{lemma}

\begin{proposition}
\label{proposition-Hartshorne-Lichtenbaum-vanishing}
\begin{reference}
\href{https://stacks.math.columbia.edu/bibliography}{Bibliography: CD (Theorem 3.1)}
\end{reference}
Let $A$ be a Noetherian local ring with completion $A^\wedge$.
Let $I \subset A$ be an ideal such that
$$
\dim V(IA^\wedge + \mathfrak p) \geq 1
$$
for every minimal prime $\mathfrak p \subset A^\wedge$ of dimension $\dim(A)$.
Then $\text{cd}(A, I) < \dim(A)$.
\end{proposition}

\begin{proof}
Since $A \to A^\wedge$ is faithfully flat we have
$H^d_{V(I)}(A) \otimes_A A^\wedge = H^d_{V(IA^\wedge)}(A^\wedge)$
by Dualizing Complexes, Lemma \href{https://stacks.math.columbia.edu/tag/0ALZ}{lemma torsion change rings (Tag 0ALZ)}.
Thus we may assume $A$ is complete.

\medskip\noindent
Assume $A$ is complete. Let $\mathfrak p_1, \ldots, \mathfrak p_n \subset A$
be the minimal primes of dimension $d$. Consider the complete local ring
$A_i = A/\mathfrak p_i$. We have $H^d_{V(I)}(A_i) = H^d_{V(IA_i)}(A_i)$
by Dualizing Complexes, Lemma
\href{https://stacks.math.columbia.edu/tag/0BJB}{lemma local cohomology and restriction (Tag 0BJB)}.
By Lemma \href{https://stacks.math.columbia.edu/tag/0EB2}{lemma cd top vanishing some module (Tag 0EB2)}
it suffices to prove the lemma for $(A_i, IA_i)$.
Thus we may assume $A$ is a complete local domain.

\medskip\noindent
Assume $A$ is a complete local domain. We can choose a prime ideal
$\mathfrak p \supset I$ with $\dim(A/\mathfrak p) = 1$.
By Lemma \href{https://stacks.math.columbia.edu/tag/0EB1}{lemma cd top vanishing (Tag 0EB1)}
it suffices to prove the lemma for $\mathfrak p$.

\medskip\noindent
By Lemma \ref{lemma-cd-bound-dualizing} it suffices to show that
$H^0_{V(\mathfrak p)}(\omega_A^\bullet) = 0$.
Recall that
$$
H^0_{V(\mathfrak p)}(\omega_A^\bullet) =
\colim \text{Ext}^0_A(A/\mathfrak p^n, \omega_A^\bullet)
$$
By Lemma \href{https://stacks.math.columbia.edu/tag/0EB5}{lemma inverse system symbolic powers (Tag 0EB5)}
we see that the colimit is the same as
$$
\colim \text{Ext}^0_A(A/\mathfrak p^{(n)}, \omega_A^\bullet)
$$
Since $\text{depth}(A/\mathfrak p^{(n)}) = 1$ we see that
these ext groups are zero by Lemma \href{https://stacks.math.columbia.edu/tag/0DWZ}{lemma sitting in degrees (Tag 0DWZ)}
as desired.
\end{proof}
\end{document}
